% =====================================================================
%  CHAPTER 4: RESULTS  (converted from Markdown)
%  Required preamble packages: graphicx, booktabs
%  ( \resizebox needs graphicx; \toprule/\midrule/\cmidrule/\bottomrule
%    need booktabs; \textsubscript is built into modern LaTeX. )
% =====================================================================

\section*{CHAPTER 4: EVALUATION METHODOLOGY}
\addcontentsline{toc}{section}{\numberline{}CHAPTER 4: EVALUATION METHODOLOGY}
\setcounter{section}{4}
\setcounter{subsection}{0}
\setcounter{figure}{0}
\setcounter{table}{0}

\subsection{Two measurements}

\textbf{Model selection} uses stratified \textit{k}-fold cross-validation over the \textbf{training split only}. Every decision in this report (preprocessing level, title weighting, truncation, hyperparameters, architecture) was made against cross-validation scores.

\textbf{The final result} uses the competition process on the \textbf{test split}, run exactly once per model when that model was final. This is the only number comparable with the published baseline.

\subsection{Metrics}

The course specifies accuracy, precision, recall, F1-score, and area under the ROC curve. All are derived from the confusion matrix. For a class $c$:

\[ \text{precision} = \frac{TP}{TP + FP}, \qquad \text{recall} = \frac{TP}{TP + FN}, \qquad F_1 = \frac{2PR}{P + R} \]

We implement these from first principles rather than importing them because the evaluation layer then carries no dependency on any modelling library and the definitions match exactly what the course covered.

\texttt{tests/test\_evaluation.py} cross-checks every metric against scikit-learn on 500 synthetic predictions -- per-class precision, recall, F1 and support; macro, micro and weighted averaging; the stratified fold splitter's proportions; and one-vs-rest AUC. \textbf{Agreement is exact to machine precision}, with a maximum deviation of $1.1 \times 10^{-16}$ on the AUC comparison.

\subsubsection{Averaging}

We report macro, micro and weighted F1 because they answer different questions. Macro averages the per-class scores unweighted, so a class the model ignores cannot be hidden. Micro pools the counts before dividing, which for a single-label multi-class problem makes it identical to accuracy -- a property we assert in a test.

The difference is easiest to see in the degenerate case. A model that always predicts \texttt{bug} on our balanced data scores (Table~\ref{tab:degenerate_case_metrics}). Accuracy rewards the third of the data it happens to get right; macro F1 reports that two of three classes score zero.

\begin{table}[H]
\centering
\caption{Evaluation metric values under a trivial majority-class prediction policy.}
\label{tab:degenerate_case_metrics}
\vspace{0.5em}
\begin{tabular}{l c}
\hline
 & value \\
\hline
accuracy & 0.3333 \\
micro F1 & 0.3333 \\
\textbf{macro F1} & \textbf{0.1667} \\
weighted F1 & 0.1667 \\
\hline
\end{tabular}
\end{table}

Because the dataset is exactly balanced (Section 2.1), macro and micro converge for any model that predicts all three classes at similar rates. We report both regardless: if resampling or filtering were introduced at any point, they would diverge immediately, and that divergence is the signal that something changed.

\textbf{Macro F1 is the figure used throughout}, matching the competition's definition.

\subsubsection{ROC and AUC}

The task is multi-class, so ROC is computed one-vs-rest: each class is taken as positive in turn with the other two pooled as negative, and the macro AUC is the unweighted mean of the three curves.

AUC is reported only for models that expose calibrated class probabilities. Models returning hard labels -- the keyword rules, \texttt{LinearSVC} -- report AUC as absent. This costs us an AUC column for two models and is the correct trade: a number computed differently from the rest of the column is worse than no number.

\subsection{Cross-validation}

Stratified \textit{k}-fold with $k = 10$ over the 1,500 training issues, seeded at 42. Each class is shuffled independently and dealt round-robin into the folds, so every fold holds the class proportions of the whole to within one item. In practice each fold contains exactly 150 issues -- 50 per class -- with 1,350 for training.

The splitter is seeded (42), so every run and every team member obtains identical folds and identical numbers.

Grid searches use $k = 5$ rather than 10, because a grid multiplies the number of fits and five folds is sufficient to \textit{rank} configurations; the selected configuration is then re-evaluated at $k = 10$.

For the neural models, cross-validation is supplemented by a \textbf{held-out validation split of 20\%} of the training data, as the course prescribes, from which training and validation loss are recorded at every epoch.

\subsection{The competition protocol}

The NLBSE'24 rules define the reported figure:

\begin{enumerate}
    \item Train a \textbf{separate classifier for each of the five repositories}.
    \item Evaluate each on that repository's 300 test issues.
    \item Score each repository as the \textbf{mean F1 over the three classes}.
    \item Report the \textbf{arithmetic mean of the five per-repository scores}.
\end{enumerate}

\subsection{Baselines}

\textbf{Floors} establish what costs no effort, so that a result can be read in context. A model reaching 0.75 means little until the reader knows that always predicting one class reaches 0.17 and hand-written keyword rules reach 0.54.

\textbf{The published baseline} is the competition's SetFit ~\cite{nlbse2024}, which is what ``how good is this?'' ultimately means in Table~\ref{tab:published_baseline_scores}

\begin{table}[h]
\centering
\caption{Published NLBSE'24 SetFit baseline F1 scores by repository.}
\label{tab:published_baseline_scores}
\vspace{0.5em}
\begin{tabular}{l c}
\hline
Repository & SetFit F1 \\
\hline
facebook/react & 0.8718 \\
tensorflow/tensorflow & 0.8644 \\
microsoft/vscode & 0.8262 \\
opencv/opencv & 0.8173 \\
bitcoin/bitcoin & 0.7555 \\
\textbf{cross-repository} & \textbf{0.8270} \\
\hline
\end{tabular}
\end{table}

\subsection{Guarding against leakage}

Three mechanisms, each with a corresponding test:

\textbf{Vectorisers are fitted inside the pipeline}, and therefore inside each fold. A test fits a pipeline on one half of a project's issues and asserts that no term occurring only in the held-out half entered the learned vocabulary -- comparing using the vectoriser's own analyser, since a naive whitespace split tokenises differently and comparing across the two would compare noise.

\textbf{Models are supplied as factories}, so no fitted state can leak between folds or between repositories.

\textbf{The error analysis re-derives its predictions} through the same protocol as the leaderboard, and a test asserts the two agree to four decimal places -- so Section 6 describes the same run Section 5 reports, not a separate one.